\clearpage
\SourceLeaf{453}
\begin{center}\large [ON THE THEORY OF THE MEASURE OF CURVATURE.]\end{center}
\noindent\hspace*{1.2em}The equation expressing the condition for
\[
pdx^2+2qdxdy+rdy^2
\]
\noindent to be the pure product of two exact differentials is as follows:
\[
\begin{aligned}
0={}&2(qq-pr)\left(\frac{\partial\partial p}{\partial y^2}
-2\frac{\partial\partial q}{\partial x\partial y}
+\frac{\partial\partial r}{\partial x^2}\right)\\
&+p\left(\frac{\partial p}{\partial y}.\frac{\partial r}{\partial y}
-2\frac{\partial q}{\partial x}.\frac{\partial r}{\partial y}
+\frac{\partial r}{\partial x}.\frac{\partial r}{\partial x}\right)\\
&+q\left(\frac{\partial p}{\partial x}.\frac{\partial r}{\partial y}
-\frac{\partial p}{\partial y}.\frac{\partial r}{\partial x}
+4\frac{\partial q}{\partial x}.\frac{\partial q}{\partial y}\right.\\
&\qquad\left.-2\frac{\partial q}{\partial x}.\frac{\partial r}{\partial x}
-2\frac{\partial p}{\partial y}.\frac{\partial q}{\partial y}\right)\\
&+r\left(\frac{\partial p}{\partial x}.\frac{\partial r}{\partial x}
-2\frac{\partial p}{\partial x}.\frac{\partial q}{\partial y}
+\frac{\partial p}{\partial y}.\frac{\partial p}{\partial y}\right).
\end{aligned}
\]

\Needspace{9\baselineskip}
\begin{center}\textbf{REMARKS.}\end{center}
\noindent\hspace*{1.2em}After \textsc{Gauss} had derived the expression for the measure of curvature of a surface, first on the assumption that the square of the line element has the form
\[
dS^2=mm(dt^2+du^2)
\]
\noindent (pp.\ 381 and 385 of this volume), and then on the assumption of geodesic coordinates for which it becomes
\[
dS^2=ds^2+mmd\theta^2
\]
\noindent (pp.\ 407 and 442 of this volume), he arrived, as the preceding note from a handbook shows, at the numerator of the expression for the measure of curvature in \textls[140]{arbitrary} coordinates by reasoning that the measure of curvature vanishes for surfaces developable onto the plane. The complete formula for the measure of curvature is then derived in the same handbook in the same way as in \S\ 11 of the \textit{Disquisitiones generales circa superficies curvas}, so it seemed unnecessary to publish these calculations. Immediately afterwards, in the same handwriting and ink, follows the note printed on pages 447 to 448, \textit{Most general solution of the problem of developing surfaces.}
\begin{flushright}\textsc{Stäckel}.\end{flushright}

